%--------------A VOTRE ATTENTION-------------%
% Les étudiants en master qui disposent de plus de 3 chapitres dans leurs travaux peuvent en complèter
% Les Membres doivent figurer dans la dernière version finale du mémoire après soutenance pour dépôt de mémoire

\documentclass{ifri}
\usepackage{titletoc}
\usepackage{tikz}
\usetikzlibrary{positioning,fit,arrows.meta}
\usepackage{tabularx}
\usepackage[super,sort&compress]{natbib} % citations en exposant, sans crochets
\hypersetup{citecolor=black}
\usepackage{booktabs}
% --- marqueurs de rédaction (à rechercher avec "À compléter" avant le dépôt) ---
\newcommand{\acompleter}[1]{\textcolor{red}{[\textbf{À compléter :} #1]}}
\newcommand{\amesurer}{\textcolor{red}{\emph{à mesurer}}}
\newcommand{\preuves}[1]{\begin{formal}\textbf{Preuves à produire.} #1\end{formal}}
\setlength{\glsdescwidth}{0.65\textwidth}
% \usepackage{lscape}

\typeMemoire{Diplôme de Licence en Informatique}
\optionFormation{Système d'Informations et Réseaux Informatiques}
\etudiant{AKPO \textbf{Sedegnan Jean Jacques}}
\titreDuMemoire{\Large\linespread{1.15}\selectfont Mise en place d'un système de détection et de prévention d'intrusion (IDS/IPS) pour la sécurisation des infrastructures réseau du Centre National Hospitalier et Universitaire Hubert Koutoukou Maga (CNHU-HKM)} %Implémention pour une meilleure sécurité dans les réseau LAN sous IPv6 // Proposition: Identification des vulnérabilités dans un reseau LAN IPv6 et mesures pour une meilleure sécurité.

\dateSoutenance{\acompleter{date de soutenance}}
%\promo{2\up{ème}}
\anneeScolaire{2025-2026}


%%maitre de mémoire
\encadrants{\acompleter{nom, prénom et fonction de l'encadrant académique IFRI et du tuteur de stage}}

%% Membres du Jury
\jurys{%
\begin{tabular}{llll}
	\acompleter{nom du président} & \acompleter{grade} & \acompleter{entité} & Président \\
	\acompleter{nom de l'examinateur} & \acompleter{grade} & \acompleter{entité} & Examinateur \\
	\acompleter{nom du rapporteur} & \acompleter{grade} & \acompleter{entité} & Rapporteur \\
\end{tabular}
}


\hypersetup{
 pdftitle={Mise en place d'un système IDS/IPS pour la sécurisation des infrastructures réseau du CNHU-HKM},
 pdfauthor={AKPO Sedegnan Jean Jacques},
 pdfsubject={Mémoire de licence, IFRI, SRI, 2025-2026},
 pdfkeywords={CNHU-HKM, IDS, IPS, Suricata, OPNsense, détection d'intrusion, sécurité réseau} 
 }

\color{bookColor}

%importation du glossaire
\loadglsentries{glossaire_reduit}

\begin{document}

\pageDeGarde
%\pageTitre

\pagecolor{white}

%% page vide
%\thispagestyle{empty}\ \clearpage


\selectlanguage{french}

% sommaire
\pagenumbering{roman}

\setcounter{tocdepth}{2}
\startlist{toc}
\printlist{toc}{}{\chapter*{Sommaire}}

%% rdedicaces
\dedicace

\acompleter{texte de la dédicace}

\newpage 

%% remerciements
\remerciements

\acompleter{remerciements personnels : direction de l'IFRI, encadrant académique, jury, famille et camarades}

Mes remerciements vont également au personnel de la Cellule des Systèmes d'Information du CNHU-HKM, et en particulier à l'équipe de l'Unité de Gestion des Réseaux Informatiques et de la Maintenance, pour son accueil, sa disponibilité et l'accès accordé à la documentation du réseau. \acompleter{nommer, avec leur fonction exacte, la cheffe de la CSI, le chef de l'UGRIM et les collègues qui vous ont accompagné ; vérifier l'orthographe de leurs noms}

\newpage 

% Résume
\resume
\selectlanguage{french}
\vspace*{-6cm}
\begin{abstract}
Ce mémoire porte sur la mise en place d'un système de détection et de prévention d'intrusion (IDS/IPS) pour sécuriser l'infrastructure réseau du Centre National Hospitalier et Universitaire Hubert Koutoukou Maga (CNHU-HKM) de Cotonou. Un audit réalisé lors d'un stage au service informatique et réseau a montré que le segment analysé ne présente ni segmentation visible ni mécanisme de détection d'intrusion. Après une étude comparative des solutions, le moteur Suricata est retenu pour un déploiement sur la plateforme de pare-feu OPNsense, prévu d'abord en mode détection puis en mode prévention, avec des règles personnalisées pour les flux DICOM et HL7. Un protocole d'évaluation dans un environnement virtualisé distinct de la production, fondé sur sept scénarios d'attaque rattachés au référentiel MITRE ATT\&CK et rejoués dans trois états (sans détection, IDS, IPS), est entièrement défini avec sa grille de mesures et les preuves à recueillir. L'expérimentation n'ayant pas été exécutée dans le cadre de ce mémoire, aucun résultat mesuré n'est présenté et aucune performance n'est affirmée sans mesure.

\paragraph{}
\textbf{Mots clés}: IDS/IPS, détection d'intrusion, Suricata, OPNsense, CNHU-HKM, sécurité réseau.
\end{abstract}

\newpage
\thispagestyle{empty}
\selectlanguage{english}
\addcontentsline{toc}{chapter}{Abstract}
\begin{abstract}

This thesis deals with the deployment of an intrusion detection and prevention system (IDS/IPS) to secure the network infrastructure of the Hubert Koutoukou Maga National Teaching Hospital (CNHU-HKM) in Cotonou. An audit carried out during an internship in the hospital's IT and network department showed that the analysed segment has neither visible segmentation nor any intrusion detection mechanism. After a comparative study of the available solutions, the Suricata engine is selected for deployment on the OPNsense firewall platform, planned first in detection mode and then in prevention mode, with custom rules for DICOM and HL7 flows. An evaluation protocol in a virtualised environment separate from production, based on seven attack scenarios mapped to the MITRE ATT\&CK framework and replayed in three states (no detection, IDS, IPS), is fully defined together with its measurement grid and the evidence to be collected. As the experiment was not carried out within this work, no measured results are reported and no performance claim is made without measurement.

\paragraph{}
\textbf{Key words}: IDS/IPS, intrusion detection, Suricata, OPNsense, CNHU-HKM, network security.
\end{abstract}

\newpage

%liste des figures
\listoffigures 
\newpage

%liste des tableaux
\listoftables
\newpage

% Les sigles et acronymes
\setglossarystyle{altlist}
\glsaddall
\printglossary[title=Liste des acronymes, toctitle=Liste des acronymes, type=\acronymtype]
\newpage

% Le glossaire proprement dit
%\setglossarystyle{super}
%\printglossary[type=main]


\pagenumbering{arabic}
\setcounter{page}{1}
%%introduction
\introduction

Le secteur de la santé connaît, depuis plusieurs années, une transformation numérique profonde : dossiers patients informatisés, imagerie médicale numérique, systèmes d'information hospitaliers interconnectés, équipements biomédicaux en réseau. Cette numérisation améliore la qualité et la rapidité de la prise en charge, mais elle expose également les établissements de santé à des risques informatiques dont les conséquences dépassent la seule perte de données : une interruption de service peut directement affecter la sécurité des patients. Les hôpitaux figurent aujourd'hui parmi les cibles privilégiées des cyberattaques, en particulier les attaques par rançongiciel, qui exploitent la criticité de leurs systèmes pour maximiser la pression sur les victimes.

Le Centre National Hospitalier et Universitaire Hubert Koutoukou Maga (CNHU-HKM) de Cotonou, établissement de référence nationale\cite{decret2015}, n'échappe pas à cette réalité. Son système d'information repose sur une infrastructure réseau qui interconnecte des services aussi divers que l'administration, les unités de soins, le laboratoire, l'imagerie médicale et les équipements biomédicaux. Un audit préliminaire mené dans le cadre de ce stage a permis de constater que le segment analysé ne présente, à ce stade, aucune segmentation visible entre ces différents usages, et qu'aucun mécanisme de détection d'intrusion n'y a été identifié : une compromission d'un poste peut donc se propager et passer inaperçue.

C'est dans ce contexte que s'inscrit le présent mémoire, réalisé à l'occasion d'un stage de \acompleter{durée du stage} au sein du service informatique et réseau du CNHU-HKM. Il traite de la mise en place d'un système de détection et de prévention d'intrusion (IDS/IPS) et propose de répondre à la problématique suivante :

\monBloc{Comment mettre en place, au CNHU-HKM, un système de détection et de prévention d'intrusion capable de repérer et de bloquer les attaques visant le réseau, tout en garantissant la disponibilité des services critiques de soins ?}

Pour y répondre, ce travail poursuit un objectif général, concevoir, déployer et valider dans un environnement de test représentatif un système IDS/IPS adapté au contexte hospitalier, décliné en cinq objectifs spécifiques :
\begin{enumerate}
\item réaliser un audit de l'existant permettant d'identifier les vulnérabilités du réseau actuel et le besoin de détection ;
\item comparer les solutions IDS/IPS disponibles et concevoir l'architecture de détection retenue ;
\item déployer Suricata, d'abord en mode détection (IDS) puis en mode prévention (IPS), sur la plateforme de pare-feu OPNsense ;
\item rédiger des règles de détection adaptées au contexte hospitalier, notamment pour les flux DICOM et HL7 ;
\item évaluer, par des scénarios d'attaque reproductibles, l'efficacité du système en comparant l'état du réseau avant et après sa mise en œuvre.
\end{enumerate}

L'intérêt de ce travail est double. Pour le CNHU-HKM, il fournit un état des lieux documenté, un système de détection testé et des recommandations concrètes de déploiement, directement exploitables par le service informatique. Pour l'auteur, il constitue une mise en application rigoureuse des compétences acquises en administration réseau et en sécurité des systèmes d'information, dans un contexte professionnel réel et à forts enjeux. La segmentation du réseau et la supervision des alertes sont traitées comme des éléments de contexte qui complètent le système, sans en être l'objet.

La démarche adoptée associe une phase d'audit, une phase de conception, une phase de réalisation en environnement de test virtualisé et une phase d'évaluation expérimentale. Le mémoire s'organise en trois chapitres. Le chapitre~\ref{chap:1} présente le cadre d'étude, le CNHU-HKM, les enjeux de la cybersécurité hospitalière et l'état de l'art de la détection et de la prévention d'intrusion. Le chapitre~\ref{chap:2} détaille la méthodologie, présente les résultats de l'audit, compare les solutions IDS/IPS et expose la conception retenue. Le chapitre~\ref{chap:3} décrit le déploiement de Suricata, les règles personnalisées, les scénarios de test, les résultats de l'évaluation avant/après et les recommandations. Le mémoire se conclut par une synthèse des apports, des limites et des perspectives de ce travail.

%\lhead[]{} \rhead[]{} \chead[]{}
\selectlanguage{french}
\fancyhead[L]{\tiny \leftmark}
\fancyhead[R]{\scriptsize \rightmark}
\fancyfoot[C]{\thepage}

\chapter{Contexte, enjeux et état de l'art de la détection d'intrusion}\label{chap:1}
 \addcontentsline{toc}{section}{Introduction}
\section*{Introduction}

Ce premier chapitre pose le cadre de l'étude. Il présente le CNHU-HKM et le déroulement du stage, puis les enjeux propres à la cybersécurité hospitalière. Il établit enfin l'état de l'art des notions et technologies mobilisées : segmentation réseau, pare-feu, détection et prévention d'intrusion, supervision de sécurité et méthodologies d'évaluation.

\section{Présentation du CNHU-HKM}

\subsection{Historique et missions}

L'établissement trouve son origine dans la loi n°~62-36 du 30 octobre 1962, qui a créé l'Hôpital de Cotonou et l'a doté de l'autonomie financière. Il a ensuite été érigé en Centre National Hospitalier par le décret n°~490/PR/MSPAS du 21 décembre 1966, puis en Centre National Hospitalier et Universitaire de Cotonou par le décret n°~73-8 du 10 janvier 1973, ses statuts étant approuvés par le décret n°~91-77 du 13 mai 1991\cite{decret2015}. Le décret n°~2012-422 du 6 novembre 2012 lui a donné son nom actuel en l'attachant à Hubert Koutoukou Maga\cite{decret2012}, et le décret n°~2015-396 du 20 juillet 2015, texte de référence sur ses attributions, son organisation et son fonctionnement, a abrogé le précédent\cite{decret2015}. Le gouvernement a par ailleurs décidé en septembre 2022 de fusionner le CNHU-HKM et l'hôpital d'instruction des armées de Cotonou\cite{fusion2022} \acompleter{préciser avec le tuteur l'état actuel de cette fusion et ses effets sur l'organisation du centre}.

Selon le décret de 2015, le CNHU-HKM est un office à caractère social et scientifique, doté de la personnalité morale et de l'autonomie financière, placé sous la tutelle du ministère en charge de la Santé et dont le siège est fixé à Cotonou, chef-lieu du département du Littoral\cite{decret2015}. Il est qualifié d'établissement de référence nationale et assure des prestations de soins préventifs, curatifs, promotionnels et réadaptatifs ; il participe en outre à la formation du personnel de santé et à la recherche en santé, et appartient à l'espace hospitalier et universitaire\cite{decret2015}. Ces trois missions, soins, enseignement et recherche, expliquent la diversité des usages que son réseau informatique doit supporter. Le centre est actuellement engagé dans un programme de réhabilitation complète de ses infrastructures, annoncé par le ministre de la Santé lors de l'installation de son nouveau conseil d'administration en octobre 2025\cite{ca2025}.

\subsection{Organisation générale}

\paragraph{Gouvernance.} Le centre est administré par un conseil d'administration dont les membres sont nommés par décret en conseil des ministres pour trois ans\cite{decret2015}. Sa composition a évolué : treize membres en 2015\cite{decret2015}, sept depuis le décret du 20 mars 2019\cite{decret2019}. Un nouveau conseil de sept membres, présidé par le représentant du ministère de la Santé, a été nommé par décret du 10 septembre 2025 pour un mandat de trois ans renouvelable\cite{ca2025nom}. Le conseil définit la politique générale de l'établissement, adopte son budget et son organigramme, et approuve le recrutement et la révocation du directeur général\cite{ca2025}.

\paragraph{Direction.} La gestion quotidienne est assurée par la Direction générale. Le directeur général est nommé par décret pris en conseil des ministres et peut être assisté d'un directeur général adjoint, nommé par arrêté du ministre de la Santé\cite{decret2015}. Le décret de 2015 place sous la coordination du directeur général adjoint \emph{la cellule informatique}, la cellule des statistiques et du contrôle de gestion et la cellule des archives et de la documentation\cite{decret2015}. Cinq structures techniques complètent la Direction générale : la Direction des Affaires Médicales, la Direction des Ressources Humaines, la Direction des Affaires Économiques, la Direction du Budget et l'Agence Comptable\cite{decret2015}. La Direction des Affaires Économiques est notamment chargée de la gestion des infrastructures, des équipements et de leur maintenance\cite{decret2015}.

\paragraph{Organes consultatifs.} Quatre comités et commissions spécialisés éclairent la gestion : le Comité de Direction, la Commission Médicale d'Établissement (CME), la Commission d'Hygiène et de Sécurité et le Comité de Lutte contre les Infections Nosocomiales\cite{decret2015}. La CME doit notamment donner son avis sur le choix du matériel médico-technique, ce qui concerne directement les équipements biomédicaux connectés au réseau.

\begin{figure}[htbp]
\centering
\small
\begin{tikzpicture}[
  box/.style={draw, rounded corners=2pt, align=center, minimum height=8mm,
              text width=2.6cm, inner sep=2pt, font=\scriptsize},
  top/.style={box, fill=gray!20, text width=4.2cm},
  node distance=6mm and 4mm]
  \node[top] (ca) {Conseil d'administration};
  \node[top, below=of ca] (dg) {Directeur général};
  \node[top, below=of dg] (dga) {Directeur général adjoint};
  \node[box, below=10mm of dga] (inf) {Cellule informatique};
  \node[box, left=of inf] (sta) {Cellules statistiques et contrôle de gestion};
  \node[box, right=of inf] (arc) {Cellule archives et documentation};
  \node[box, below=14mm of inf, xshift=-4.1cm] (dam) {Direction des Affaires Médicales};
  \node[box, right=of dam] (drh) {Direction des Ressources Humaines};
  \node[box, right=of drh] (dae) {Direction des Affaires Économiques};
  \node[box, right=of dae] (dbu) {Direction du Budget};
  \node[draw, dashed, inner sep=4pt, fit=(dam)(drh)(dae)(dbu), label={[font=\scriptsize]below:Structures techniques de la Direction générale}] (st) {};
  \draw (ca) -- (dg) -- (dga);
  \draw (dga) -- ++(0,-0.5) -| (sta);
  \draw (dga) -- (inf);
  \draw (dga) -- ++(0,-0.5) -| (arc);
\end{tikzpicture}
\caption{Organisation simplifiée du CNHU-HKM, d'après le décret n°~2015-396 (l'Agence Comptable et les comités spécialisés ne sont pas représentés). À valider avec le service informatique.}
\label{fig:organigramme}
\end{figure}

\subsection{Le service informatique et réseau}

Le décret de 2015 ne définit pas de « service informatique » en tant que direction : il prévoit une \emph{cellule informatique} coordonnée par le directeur général adjoint\cite{decret2015}. Dans la pratique, l'informatique de l'établissement est portée par la Cellule des Systèmes d'Information (CSI), chargée de l'infrastructure informatique et des télécommunications, et qui intervient dans les services médicaux, administratifs et techniques\cite{kaho2026}. Cette cellule correspond vraisemblablement à la cellule informatique du décret \acompleter{à confirmer avec le tuteur}. Elle se compose de trois unités (figure~\ref{fig:csi}) : l'Unité de Gestion des Réseaux Informatiques et de la Maintenance (UGRIM), l'Unité d'Exploitation et de Développement Informatique (UEDI) et l'Unité Administrative des Archives et Médicales (UAAM)\cite{kaho2026}.

\begin{figure}[htbp]
\centering
\small
\begin{tikzpicture}[
  box/.style={draw, rounded corners=2pt, align=center, minimum height=9mm, text width=4.2cm, font=\scriptsize},
  node distance=8mm]
  \node[box, fill=gray!20] (csi) {Cellule des Systèmes d'Information (CSI)};
  \node[box, below=of csi, xshift=-5cm] (ugrim) {UGRIM\\ réseaux et maintenance};
  \node[box, below=of csi] (uedi) {UEDI\\ exploitation et développement};
  \node[box, below=of csi, xshift=5cm] (uaam) {UAAM\\ archives administratives et médicales};
  \draw (csi) -- ++(0,-0.6) -| (ugrim);
  \draw (csi) -- (uedi);
  \draw (csi) -- ++(0,-0.6) -| (uaam);
\end{tikzpicture}
\caption{Organisation de la CSI, d'après le rapport de stage\cite{kaho2026}. À valider avec le tuteur.}
\label{fig:csi}
\end{figure}

Le stage s'est déroulé à l'UGRIM. Cette unité gère et fait évoluer les infrastructures réseau et télécoms de l'hôpital : réseaux filaires, Wi-Fi, téléphonie et installations audiovisuelles. Elle agit sur le plan préventif (installation, configuration, sécurisation) et curatif (dépannage, résolution d'incidents), et reçoit directement les requêtes des utilisateurs\cite{kaho2026}. Elle déploie, supervise et maintient les commutateurs, routeurs et points d'accès, et met en place des procédures pour protéger les connexions et prévenir les intrusions\cite{kaho2026}. L'UEDI s'occupe des applications métiers, de l'administration des serveurs physiques et virtuels, du stockage central et des sauvegardes ; l'UAAM de la gestion des archives médicales et administratives, y compris leur numérisation\cite{kaho2026}. L'effectif de chaque unité et le parc géré sont \acompleter{à compléter avec le tuteur}. L'audit du 9~octobre a déjà révélé 85 hôtes actifs sur un seul segment /23, ce qui constitue un premier indicateur de la taille du parc.

\section{Déroulement du stage}

\subsection{Accueil et intégration}
Le stage a été effectué au sein de l'UGRIM, rattachée à la CSI. L'encadrement est assuré par le chef de l'UGRIM et l'équipe technique de l'unité\cite{kaho2026}. \acompleter{date de début, noms de vos encadrants et modalités d'intégration, propres à votre stage}

\subsection{Tâches confiées}
Les stagiaires de l'unité sont associés à des missions courantes : maintenance des postes de travail, support aux utilisateurs, installation d'imprimantes réseau, câblage et sertissage de câbles Ethernet, configuration de commutateurs, installation d'équipements (points d'accès Wi-Fi, caméras, onduleurs) et assistance aux travaux de segmentation par VLAN\cite{kaho2026}. Ces activités donnent une vision concrète de l'infrastructure sur laquelle porte ce mémoire. \acompleter{ajuster à vos propres tâches et à leur part dans le service}

\subsection{Genèse et validation du sujet}

Le sujet de ce mémoire a été défini en concertation avec le tuteur de stage, à partir d'un double constat : d'une part, l'absence d'une politique de segmentation réseau formalisée au sein de l'infrastructure existante ; d'autre part, la volonté de renforcer les compétences et les moyens en cybersécurité de la cellule, avec en particulier un pare-feu avancé associé à des systèmes de détection et de prévention d'intrusion (IDS/IPS), axe de modernisation déjà identifié au sein de l'unité\cite{kaho2026} \acompleter{confirmer avec le tuteur le besoin qu'il a exprimé}. Ce travail s'inscrit en complément des travaux de segmentation par VLAN menés par l'UGRIM, en ajoutant le filtrage entre zones, la détection d'intrusion et la supervision. Le périmètre du projet, les autorisations d'accès à la documentation réseau et les modalités de travail ont fait l'objet d'un accord formalisé avec le service, présenté en annexe~\ref{annexe1}.

\section{Enjeux de la cybersécurité en milieu hospitalier}

\subsection{Spécificités d'un système d'information de santé}

Un système d'information hospitalier (SIH) se distingue des systèmes d'information classiques par trois exigences qui entrent parfois en tension. La disponibilité y est critique : une interruption, même brève, d'un service comme le dossier patient informatisé, l'imagerie ou la prescription électronique peut retarder une prise en charge. La confidentialité des données de santé, protégées par le secret médical et par les cadres légaux de protection des données personnelles, impose des contrôles d'accès stricts. Enfin, l'hétérogénéité du parc, avec des postes bureautiques, des serveurs applicatifs, mais aussi des équipements biomédicaux souvent anciens, rarement conçus avec la sécurité comme exigence première et difficiles à mettre à jour, complique l'application de mesures de sécurité uniformes.

\subsection{Panorama des menaces visant les hôpitaux}

Le premier rapport de l'agence européenne ENISA consacré au secteur de la santé couvre la période de janvier 2021 à mars 2023 et analyse 215 incidents rendus publics dans l'Union européenne et les pays voisins\cite{enisa2023,enisa2023slides}. Les prestataires de soins y représentent 53~\% des incidents et les hôpitaux à eux seuls 42~\%\cite{enisa2023}. Le rançongiciel est l'une des principales menaces du secteur (54~\% des incidents), notamment du fait de la valeur des données des patients, qui sont l'actif le plus ciblé par ces attaques (30~\%)\cite{enisa2023press,infosec2023}. Le rapport observe aussi que seules 27~\% des organisations interrogées disposent d'un programme dédié de défense contre les rançongiciels\cite{heimdal2023}, et retient parmi les défis du secteur les vulnérabilités des dispositifs médicaux\cite{enisa2023slides}. Les conséquences dépassent la perte de données : ces attaques ont provoqué des interruptions de services de soins et ont été associées à des décès de patients\cite{infosec2023}.

Ces constats concernent l'Europe et ne décrivent pas la situation du Bénin ; ils sont retenus ici comme référence sur la nature des menaces. Ils appuient la justification de la segmentation, qui vise précisément à empêcher qu'un poste compromis, point d'entrée fréquent d'un rançongiciel, ne permette d'atteindre les serveurs applicatifs.

\subsection{Cadre réglementaire applicable}

Au Bénin, la protection des données à caractère personnel relève du Livre~V du Code du numérique, institué par la loi n°~2017-20 du 20 avril 2018 et modifiée par la loi n°~2020-35 du 6 janvier 2021\cite{codenum,asin}. Le code identifie l'Autorité de Protection des Données à caractère Personnel (APDP) comme autorité de contrôle\cite{codenum}. Tout traitement de données personnelles y est soumis à une autorisation préalable de l'APDP, conformément à l'article~405\cite{apdp405}. Les articles 425 et 426, relatifs à la gestion et à la sécurité des données, sont ceux auxquels l'APDP renvoie pour la protection des données de santé\cite{apdp2024}.

Interrogée sur les données de santé, une représentante de l'APDP recommande, sans préjudice des règles classiques de sécurité informatique, d'utiliser l'authentification à deux facteurs, de bâtir des réseaux et systèmes sécurisés, de chiffrer les données sensibles, de limiter les accès, de réaliser des sauvegardes régulières et de sensibiliser les employés\cite{apdp2024}. Ces recommandations recoupent les mesures du présent mémoire : segmentation, filtrage, supervision et durcissement des accès d'administration. Le Code du numérique prévoit en outre une agence nationale chargée de contrôler le niveau de sécurité des systèmes d'information, l'ANSSI-Bénin\cite{codenum}.

Les référentiels internationaux ISO/IEC 27001, ISO/IEC 27799 et PGSSI-S sont mobilisés comme cadre de référence méthodologique et non comme obligation réglementaire locale \acompleter{vérifier chaque référence avant citation}.

\section{Place de l'IDS/IPS dans l'architecture de sécurité}

La sécurité d'un réseau repose sur la superposition de plusieurs mesures complémentaires, la défense en profondeur. Trois d'entre elles encadrent le système étudié dans ce mémoire.

La segmentation réseau divise une infrastructure en zones séparées, de façon à ce que la compromission d'un équipement ne se propage pas automatiquement aux autres. Elle s'appuie sur des réseaux virtuels (VLAN, \emph{Virtual Local Area Network}) et sur des zones de sécurité telles que la zone démilitarisée (DMZ). En milieu hospitalier, elle vise à isoler les équipements biomédicaux, dont le micrologiciel est rarement mis à jour, des postes administratifs et des serveurs. Les flux métiers propres au secteur, en particulier le protocole DICOM (\emph{Digital Imaging and Communications in Medicine}) pour l'imagerie et le standard HL7 (\emph{Health Level 7}) pour l'échange de données cliniques, doivent être explicitement identifiés.

Le pare-feu contrôle le passage entre zones selon une politique explicite. Un pare-feu à état (\emph{stateful}) conserve la table des connexions actives et autorise le trafic de retour légitime, là où un pare-feu sans état traite chaque paquet isolément\cite{nist80041}. Il décide toutefois à partir d'adresses, de ports et de protocoles, et n'analyse pas nécessairement le comportement d'un flux pourtant autorisé. OPNsense est une distribution de pare-feu open source fondée sur FreeBSD, dérivée du projet pfSense, qui propose un filtrage à état, une interface web d'administration, un VPN et un service de détection et de prévention d'intrusion fondé sur Suricata\cite{opnsensedoc}. Elle a fait l'objet d'une évaluation de sécurité par le laboratoire jtsec\cite{jtsec} \acompleter{préciser la nature de l'évaluation (certification, qualification) d'après le rapport}, ce qui constitue un élément objectif en faveur de sa fiabilité.

L'IDS/IPS comble cette limite : il inspecte le contenu et le comportement du trafic, y compris celui que le pare-feu autorise, et signale ou bloque ce qui ressemble à une attaque. Il constitue ainsi la brique centrale de ce travail.

\section{Détection et prévention d'intrusion}

\subsection{Principes}

Un système de détection d'intrusion (IDS, \emph{Intrusion Detection System}) analyse le trafic réseau ou l'activité d'un système à la recherche de signes d'activité malveillante et génère une alerte lorsqu'il en détecte, sans intervenir sur le trafic. Un système de prévention d'intrusion (IPS, \emph{Intrusion Prevention System}) va plus loin : placé en coupure sur le flux, il peut bloquer activement le trafic identifié comme malveillant\cite{nist80094}. On distingue les systèmes réseau (NIDS), qui observent le trafic en un point du réseau, et les systèmes hôte (HIDS), qui surveillent l'activité d'une machine (journaux, fichiers, processus).

La détection s'appuie sur deux approches complémentaires. La détection par signatures compare le trafic à une base de motifs d'attaques connues : elle est précise mais limitée aux menaces déjà répertoriées. La détection par anomalies repère les écarts par rapport à un comportement de référence : elle peut déceler des menaces inédites, au prix d'un taux de faux positifs plus élevé\cite{nist80094}. Un faux positif est une alerte déclenchée par un trafic légitime ; un faux négatif est une attaque qui passe sans alerte. Pour un IPS, un faux positif est plus grave encore, puisqu'il peut bloquer un flux de soins légitime.

\subsection{Positionnement d'une sonde}

Le mode de déploiement conditionne ce que le système peut faire. En mode passif, la sonde reçoit une copie du trafic (port miroir ou TAP) : elle détecte et alerte sans jamais perturber le flux, ce qui en fait le mode adapté à une première phase. En mode en coupure (\emph{inline}), tout le trafic traverse la sonde, qui peut le bloquer, mais qui devient aussi un point de défaillance : une panne ou une saturation interrompt la circulation\cite{nist80094}. Cette contrainte est déterminante en milieu hospitalier et motive une mise en service progressive, de l'IDS vers l'IPS.

\subsection{Suricata}

Suricata est un moteur de détection et de prévention d'intrusion open source, capable de fonctionner en mode IDS comme en mode IPS. Il analyse le trafic à l'aide de règles, dont le jeu ET Open (\emph{Emerging Threats Open}), mis à jour par une communauté de chercheurs en sécurité, et accepte des règles personnalisées. Une règle comprend une action (\texttt{alert} pour signaler, \texttt{drop} pour bloquer en mode IPS), un en-tête décrivant le protocole, les adresses, les ports et le sens du flux, puis des options entre parenthèses (message, contenu recherché, seuils, identifiant \texttt{sid}). Suricata journalise les événements dans un format structuré, \texttt{eve.json}, exploitable par un outil de supervision tiers\cite{suricatadoc}. La documentation de la version 8.0.7 présente par ailleurs un mode pare-feu expérimental, dont la politique par défaut est le rejet\cite{suricata} ; il n'est pas retenu ici et reste une piste d'évolution.

\subsection{Limites}

Un IDS/IPS ne remplace pas les autres mesures. Il ne voit pas le contenu du trafic chiffré, il demande un réglage continu des règles pour limiter les faux positifs, et son efficacité dépend de la qualité et de la mise à jour de son jeu de signatures. Sa charge de traitement augmente avec le débit inspecté, ce qui impose de mesurer son impact sur les performances.

\section{Exploitation des alertes}

Les alertes d'une sonde ne servent que si quelqu'un les lit. Un système de gestion des informations et des événements de sécurité (SIEM, \emph{Security Information and Event Management}) centralise les journaux de plusieurs équipements, les corrèle et alerte sur les comportements suspects ; certains offrent une réponse active, comme le blocage temporaire d'une adresse source. Wazuh est une solution open source structurée autour d'un gestionnaire, d'un indexeur et d'un tableau de bord web, qui peut collecter les journaux d'un pare-feu ou d'une sonde Suricata\cite{wazuhdoc}. Dans ce mémoire, Wazuh est utilisé comme outil complémentaire pour visualiser les alertes de Suricata ; il n'en est pas l'objet principal.

\section{Méthodologies d'évaluation}

MITRE ATT\&CK est une base de connaissances publique qui répertorie les tactiques et techniques employées par les attaquants, organisées selon les phases d'une attaque (reconnaissance, accès initial, mouvement latéral, etc.)\cite{attck}. Elle permet de nommer chaque scénario de test selon une nomenclature reconnue.

L'évaluation consiste à rejouer les mêmes scénarios d'attaque dans un environnement de test identique, successivement sans système de détection (état « avant »), puis avec Suricata en IDS et en IPS (états « après »). Cette comparaison mesure l'apport du système à travers le nombre d'attaques détectées ou bloquées, le délai de détection, le nombre de faux positifs et l'impact sur les performances. Les limites d'un environnement virtualisé sont discutées en section~\ref{sec:limites}.

\section{Travaux et solutions existants}

\acompleter{identifier deux à quatre mémoires, articles ou études de cas portant sur la sécurisation de réseaux hospitaliers ou sur l'usage combiné de pare-feu open source, d'IDS/IPS et de SIEM. Pour chacun : approche, résultats, puis positionnement de ce mémoire (contexte local à ressources limitées, parc hétérogène, évaluation avant/après). Aucune référence n'est ajoutée ici tant qu'elle n'a pas été lue et vérifiée.}

\addcontentsline{toc}{section}{Conclusion}
\section*{Conclusion}

Ce chapitre a situé le cadre institutionnel du stage et mis en évidence les exigences propres à la sécurité d'un système d'information hospitalier : disponibilité, confidentialité et hétérogénéité du parc. Il a présenté les fondements théoriques mobilisés : segmentation, filtrage par pare-feu, détection et prévention d'intrusion, supervision des événements de sécurité, ainsi que le référentiel MITRE ATT\&CK et la démarche d'évaluation comparative. Le chapitre suivant présente l'audit du réseau existant, compare les solutions candidates et expose la conception de l'architecture retenue.

 
\chapter{Audit de l'existant, étude comparative des solutions IDS/IPS et conception}\label{chap:2}
 \addcontentsline{toc}{section}{Introduction}
\section*{Introduction}

Ce chapitre décrit la démarche suivie, expose la méthode du futur audit du réseau du CNHU-HKM et les éléments d'architecture documentés par le service, compare les solutions IDS/IPS candidates, puis expose la conception du système de détection retenu. Les éléments de conception de la section~\ref{sec:conception} sont des propositions de travail, à valider avec le tuteur avant leur mise en œuvre.

\section{Méthodologie}

\subsection{Étapes du projet}

Le projet s'organise en cinq phases successives, étalées sur \acompleter{X} semaines : un audit du réseau existant, une phase de conception de l'architecture cible, une phase de réalisation en environnement de test virtualisé, une phase d'évaluation expérimentale par des scénarios d'attaque, et une phase de synthèse donnant lieu à des recommandations pour un déploiement progressif au CNHU-HKM. La phase de conception (phase 2) est réalisée et documentée dans ce mémoire. La phase d'audit (phase 1) doit être conduite selon la méthode décrite dans la suite de ce chapitre ; ses résultats ne sont pas présentés ici. Les phases 3 et 4 n'ont pas été exécutées : elles font l'objet, au chapitre~\ref{chap:3}, d'un protocole entièrement spécifié et prêt à mettre en œuvre. Le tableau~\ref{tab:calendrier} fixe le calendrier prévisionnel de ces phases ; les dates et durées effectives doivent être reportées depuis le journal de bord, aucune durée n'étant indiquée tant qu'elle n'est pas attestée.

\begin{table}[htbp]
\centering
\small
\begin{tabularx}{\textwidth}{cllX}
\hline
\textbf{Phase} & \textbf{Semaines effectives} & \textbf{Contenu} & \textbf{Livrable} \\
\hline
1 & \acompleter{dates} & Audit du réseau existant, découverte d'hôtes et de services & Rapport d'audit et constats \\
2 & \acompleter{dates} & Conception de l'architecture cible, zones et flux & Plan d'adressage, matrice de flux \\
3 & \acompleter{dates} & Réalisation du laboratoire virtualisé (OPNsense, Suricata, Wazuh) & Plateforme de test documentée \\
4 & \acompleter{dates} & Évaluation expérimentale : scénarios S1 à S7, états A, B et C & Tableau de résultats et preuves \\
5 & \acompleter{dates} & Synthèse, recommandations de déploiement progressif & Rapport final \\
\hline
\end{tabularx}
\caption{Calendrier prévisionnel de réalisation ; les semaines effectives sont à reporter depuis le journal de bord}
\label{tab:calendrier}
\end{table}

\subsection{Cadre éthique et autorisations}

Les manipulations du futur audit réseau (découverte d'hôtes et balayage de ports) devront être menées dans le respect d'un cadre formalisé avec le service informatique et réseau du CNHU-HKM. Le périmètre autorisé pour la consultation de la documentation réseau et pour les opérations de découverte sera fixé par une autorisation écrite du tuteur de stage \acompleter{joindre le document signé en annexe~H}, qui rappellera également l'interdiction de tout test d'intrusion actif sur le réseau de production sans accord préalable explicite. Ce cadre impose qu'aucune donnée patient ne soit consultée ni manipulée. Les informations sensibles relatives à l'infrastructure (adresses IP réelles, noms de machines) seront anonymisées lorsque leur divulgation n'apportera pas de valeur pédagogique. Les scénarios d'attaque du chapitre~\ref{chap:3} sont, quant à eux, exclusivement prévus dans le laboratoire virtualisé : ils ne doivent pas être exécutés sur le réseau de production.

\section{Audit du réseau : méthode envisagée et éléments documentés}

\subsection{Méthode de découverte envisagée}

L'audit réseau sera mené à l'aide de l'outil Nmap, en deux temps ; la version exacte de l'outil sera précisée lors de l'exécution. La première phase sera une découverte d'hôtes actifs par requête ARP (\texttt{nmap -sn}), sans interrogation de services, sur le segment de données du service informatique. La seconde consistera en un balayage ciblé sur sept ports associés à des services web et de vidéosurveillance (80, 443, 554, 8000, 8080, 8899, 37777), avec détection de version des services (\texttt{nmap -sV}), afin d'identifier la nature des équipements exposant un service sur ces ports. Aucun test d'intrusion, de force brute ni capture de trafic ne sera réalisé : la découverte se limitera aux services exposés, conformément au cadre éthique défini plus haut. Les sorties complètes des deux phases seront archivées sans modification en annexe~E.

\subsection{Nouvel audit programmé}

Les résultats d'un premier audit réalisé pendant le stage ne sont pas repris dans ce mémoire : un nouvel audit réseau doit être conduit ultérieurement, selon la méthode ci-dessus, pour établir un état des lieux à jour. Aucun résultat, aucune adresse IP, aucun port détecté ni aucune interprétation n'est donc présenté par anticipation : les constats, le tableau de synthèse par constructeur d'équipement et la liste des services exposés seront rédigés exclusivement à partir des sorties effectivement recueillies. \acompleter{date et périmètre exacts du nouvel audit, à fixer avec le tuteur ; tableau de synthèse des constructeurs et nombre d'hôtes, à extraire des futures sorties Nmap}

\subsection{Architecture physique du réseau}

Le service a fourni un schéma de l'architecture du réseau du CNHU-HKM\cite{archi}, dont la date exacte reste à préciser. Il montre une topologie hiérarchique : un routeur Cisco relié à Internet, un switch fédérateur (Dell) qui sert de cœur, puis des liaisons en fibre optique vers des armoires de brassage et des switches d'accès répartis par service (urgences, soins, imagerie, laboratoire, banque de sang, administration, pharmacies, etc.). Les postes et des points d'accès Wi-Fi s'y raccordent, et un serveur applicatif est relié au cœur (figure~\ref{fig:existant}). Le schéma distingue par couleur les switches Cisco existants de ceux à remplacer, par un Cisco simple ou par un Cisco à port SFP ; la majorité des switches d'accès y est marquée à remplacer.

\begin{figure}[htbp]
\centering
\small
\begin{tikzpicture}[
  box/.style={draw, rounded corners=2pt, align=center, minimum height=8mm, text width=3.4cm, font=\scriptsize},
  node distance=7mm and 6mm]
  \node[box] (net) {Internet};
  \node[box, right=of net] (rtr) {Routeur Cisco};
  \node[box, right=of rtr, fill=gray!25] (core) {Switch fédérateur (Dell)};
  \node[box, below=14mm of rtr] (srv) {Serveur applicatif};
  \node[box, below=14mm of core] (ex) {Armoires avec switches Cisco existants};
  \node[box, right=of ex] (rem) {Switches d'accès à remplacer (services)};
  \node[box, below=10mm of ex] (ap) {Points d'accès Wi-Fi};
  \node[box, dashed, below=14mm of srv] (sonde) {Sonde IDS (port miroir)\\ \tiny emplacement proposé};
  \draw (net) -- (rtr) -- (core);
  \draw (core) -- node[left, font=\tiny]{FO} (srv);
  \draw (core) -- node[left, font=\tiny]{FO} (ex);
  \draw (core.south east) -- ++(0,-0.4) -| node[above, pos=0.2, font=\tiny]{FO} (rem);
  \draw (ex) -- (ap);
  \draw[dashed] (core.south west) -- ++(-0.9,-0.5) |- (sonde.east);
\end{tikzpicture}
\caption{Architecture physique simplifiée, redessinée d'après le schéma du service\cite{archi}. L'emplacement de la sonde est une proposition}
\label{fig:existant}
\end{figure}

Quatre enseignements se dégagent pour ce travail.
\begin{itemize}
\item \textbf{Un point de convergence unique.} Le trafic entre services transite par le switch fédérateur et le routeur. C'est un emplacement d'observation idéal : un port miroir sur le fédérateur donne à une sonde passive la vue sur l'essentiel des flux, sans effet sur le trafic. C'est aussi un point de défaillance unique, ce qui pèse sur le choix d'un mode en coupure.
\item \textbf{Aucun pare-feu, VLAN ni adressage n'apparaît sur le schéma.} Cette absence de représentation ne préjuge pas de l'existence de ces éléments : le nouvel audit et les échanges avec le service devront lever cette ambiguïté.
\item \textbf{Un parc de switches hétérogène et en renouvellement.} Les capacités (port miroir, VLAN, contrôle ARP) dépendent du modèle, ce qui conditionne la faisabilité de la sonde et de la segmentation. \acompleter{modèle exact du switch fédérateur Dell et prise en charge du port miroir}
\item \textbf{Une limite de visibilité.} Le trafic entre deux postes d'un même switch d'accès ne traverse pas le cœur et échappe à une sonde placée à cet endroit.
\end{itemize}
La date exacte du schéma restant à préciser, il doit être validé par le service avant d'être exploité comme référence ; le nouvel audit permettra notamment de vérifier le nombre et le positionnement des points d'accès Wi-Fi.

\subsection{Inventaire, flux et usages}

\acompleter{inventaire détaillé par zone fonctionnelle, type d'équipement et criticité estimée, à établir après le nouvel audit en croisant ses sorties avec la documentation interne et les échanges avec le personnel technique}

\acompleter{flux métiers connus du service informatique (accès au SIH, à l'imagerie, à la messagerie), décrits à partir des échanges avec le personnel, même sans capture de trafic}

\section{Analyse de risques simplifiée}

L'analyse s'inspire d'une version simplifiée de la méthode EBIOS Risk Manager : identification des biens essentiels, des scénarios redoutés, puis cotation de la vraisemblance et de la gravité sur une échelle de 1 (faible) à 4 (critique). Le tableau~\ref{tab:risques} présente des scénarios de risque anticipés à partir des menaces documentées au chapitre~\ref{chap:1} et de l'architecture schématisée plus haut ; ils devront être confirmés ou ajustés à la lumière du nouvel audit. Les cotations sont une proposition de l'auteur, \acompleter{à valider avec le tuteur}.

\begin{table}[htbp]
\centering
\small
\begin{tabularx}{\textwidth}{lXccc}
\hline
\textbf{Réf.} & \textbf{Scénario de risque} & \textbf{Vrais.} & \textbf{Gravité} & \textbf{Niveau} \\
\hline
R1 & Propagation d'un rançongiciel depuis un poste administratif vers les serveurs applicatifs, faute de segmentation & 3 & 4 & 12 \\
R2 & Compromission d'un équipement biomédical, faute d'authentification robuste et d'isolement & 2 & 4 & 8 \\
R3 & Intrusion non détectée, faute de supervision centralisée & 3 & 3 & 9 \\
R4 & Accès d'un visiteur ou d'un équipement non maîtrisé au réseau interne & 2 & 3 & 6 \\
\hline
\end{tabularx}
\caption{Scénarios de risque (niveau = vraisemblance $\times$ gravité)}
\label{tab:risques}
\end{table}

La hiérarchisation place la propagation latérale (R1) en tête, suivie de l'absence de détection (R3). Un IDS/IPS répond directement à ces deux risques : il repère les comportements de propagation et comble l'absence de détection. Il justifie le choix du système étudié dans ce mémoire.

\section{Étude comparative des solutions IDS/IPS}

\subsection{Critères}

Les solutions ont été comparées selon des critères issus des contraintes du contexte : coût de licence nul, prise en charge des modes IDS et IPS, qualité et disponibilité des règles de détection, intégration à la plateforme de pare-feu, format de journalisation exploitable, maturité et documentation, et charge de déploiement sur un laboratoire aux ressources limitées.

\subsection{Moteurs de détection}

\begin{table}[htbp]
\centering
\small
\begin{tabularx}{\textwidth}{lXXX}
\hline
\textbf{Critère} & \textbf{Suricata} & \textbf{Snort} & \textbf{Zeek} \\
\hline
Nature & moteur IDS/IPS à règles & moteur IDS/IPS à règles & analyseur de trafic et de comportements \\
Modes IDS et IPS & oui\cite{suricatadoc} & oui & surtout analyse, pas conçu comme IPS \\
Règles & ET Open et règles locales & règles communautaires et locales & scripts d'analyse \\
Intégration à OPNsense & native\cite{opnsensedoc} & non native & non \\
Journal structuré & \texttt{eve.json} & formats propres & journaux structurés \\
\hline
\end{tabularx}
\caption{Comparaison des moteurs de détection (appréciations de l'auteur)}
\label{tab:comparatif}
\end{table}

Comme le synthétise le tableau~\ref{tab:comparatif}, Suricata est retenu : il offre les modes IDS et IPS, un large jeu de règles ouvert, un journal structuré facile à exploiter et surtout une intégration native à OPNsense, ce qui limite la charge de déploiement. Zeek, plus orienté vers l'analyse du trafic, est un complément possible plutôt qu'un substitut.

\subsection{Modes de déploiement}

Deux modes de déploiement ont été envisagés. Le premier intègre Suricata au pare-feu OPNsense, qui devient à la fois point de filtrage et point d'inspection ; il est simple à mettre en œuvre et à administrer. Le second place une sonde dédiée, alimentée par un port miroir, qui ne perturbe jamais le trafic mais ne peut pas bloquer. Le premier est retenu pour ce travail, car il permet de passer de l'IDS à l'IPS sans changer d'architecture. Le second reste adapté à une phase d'observation sur le réseau de production.

\subsection{Pare-feu et supervision}

OPNsense est retenu comme plateforme hôte pour son caractère open source, son filtrage à état et son intégration de Suricata\cite{opnsensedoc}. Wazuh est retenu comme outil complémentaire de visualisation des alertes, pour ses agents, sa corrélation et sa réponse active dans un seul produit open source\cite{wazuhdoc}.

\section{Conception du système de détection}\label{sec:conception}

\subsection{Zones surveillées et plan d'adressage}

Le système inspecte le trafic qui circule entre des zones fonctionnelles. À partir de la segmentation hospitalière décrite au chapitre~\ref{chap:1} et de l'architecture documentée, cinq zones sont retenues, chacune portée par un VLAN ; leur calage définitif sur les réseaux réels interviendra après le nouvel audit. Le plan d'adressage ci-dessous est celui du laboratoire, distinct du plan d'adressage de production, qui sera documenté à l'issue du nouvel audit. Le pare-feu en est la passerelle (adresse .1 de chaque réseau).

\begin{table}[htbp]
\centering
\small
\begin{tabular}{llll}
\hline
\textbf{VLAN} & \textbf{Zone} & \textbf{Réseau} & \textbf{Contenu type} \\
\hline
10 & Administration & 10.10.10.0/24 & postes administratifs, poste d'administration \\
20 & Soins & 10.10.20.0/24 & postes des unités de soins \\
30 & Biomédical & 10.10.30.0/24 & équipements biomédicaux, imagerie \\
40 & Serveurs & 10.10.40.0/24 & serveurs applicatifs, annuaire, serveur Wazuh \\
50 & Invités & 10.10.50.0/24 & visiteurs, accès Internet seul \\
\hline
\end{tabular}
\caption{Zones surveillées et plan d'adressage du laboratoire (proposition)}
\label{tab:vlan}
\end{table}

\subsection{Positionnement de la sonde}

Dans l'architecture cible, Suricata est placé sur le pare-feu OPNsense, par lequel transitent tous les flux entre zones et vers l'extérieur (figure~\ref{fig:cible}). Il inspecte ainsi les flux inter-zones, les plus révélateurs d'une propagation latérale, ainsi que le trafic entrant et sortant.

Sur le réseau réel, le schéma de l'architecture fait ressortir deux emplacements possibles. Le premier est un port miroir du switch fédérateur, qui alimente une sonde passive : il ne perturbe pas le trafic et convient à la phase IDS. Le second est l'insertion de la sonde en coupure entre le routeur et le switch fédérateur : il permet le blocage, mais fait de la sonde un point de défaillance et impose un mécanisme de contournement ou une redondance. Le laboratoire reproduit le second cas, car il permet de tester l'IDS comme l'IPS.

\begin{figure}[htbp]
\centering
\small
\begin{tikzpicture}[
  zone/.style={draw, rounded corners=2pt, align=center, minimum height=9mm, text width=2.2cm, font=\scriptsize},
  node distance=3mm]
  \node[zone, fill=gray!25, text width=4.2cm] (fw) {Pare-feu OPNsense\\ \tiny Suricata : sonde IDS puis IPS};
  \node[zone, above=8mm of fw, text width=3.6cm] (wan) {Réseau externe / Internet};
  \node[zone, below=14mm of fw, xshift=-5.4cm] (adm) {VLAN 10\\Administration};
  \node[zone, right=of adm] (soi) {VLAN 20\\Soins};
  \node[zone, right=of soi] (bio) {VLAN 30\\Biomédical};
  \node[zone, right=of bio] (srv) {VLAN 40\\Serveurs\\ \tiny + Wazuh};
  \node[zone, right=of srv] (inv) {VLAN 50\\Invités};
  \draw (wan) -- (fw);
  \foreach \z in {adm,soi,bio,srv,inv} \draw (fw) -- (\z);
\end{tikzpicture}
\caption{Architecture cible : Suricata inspecte les flux entre zones sur le pare-feu. À remplacer par un schéma draw.io si souhaité}
\label{fig:cible}
\end{figure}

\subsection{Flux de référence}

Pour écrire des règles pertinentes, il faut connaître les flux attendus. Le tableau~\ref{tab:flux} liste les flux légitimes entre zones ; ils servent de base aux règles de détection, car tout flux hors de cette liste est suspect. Les ports applicatifs du SIH et des équipements biomédicaux sont \acompleter{à confirmer auprès du service}. Le même tableau fonde la politique de filtrage du pare-feu, qui applique le principe « tout interdit sauf exception »\cite{nist80041}.

\begin{table}[htbp]
\centering
\small
\begin{tabularx}{\textwidth}{llXl}
\hline
\textbf{Source} & \textbf{Destination} & \textbf{Service} & \textbf{Décision} \\
\hline
Administration & Serveurs & HTTPS 443, applications du SIH & Autoriser \\
Soins & Serveurs & applications du SIH, DICOM 104, HL7 (MLLP) & Autoriser \\
Biomédical & Serveurs (imagerie) & DICOM uniquement & Autoriser \\
Biomédical & Autres zones & tout & Interdire \\
Administration & Toutes zones & SSH 22, RDP 3389 depuis les postes d'administration & Autoriser, journaliser \\
Toutes zones & Serveurs (annuaire) & DNS 53, DHCP 67/68 & Autoriser \\
Hôtes supervisés & Serveur Wazuh & agent 1514, enregistrement 1515 & Autoriser \\
Pare-feu & Serveur Wazuh & syslog 514 & Autoriser \\
Invités & Internet & DNS, HTTP, HTTPS & Autoriser \\
Invités & Zones internes & tout & Interdire \\
\hline
\end{tabularx}
\caption{Flux de référence entre zones (proposition)}
\label{tab:flux}
\end{table}

\subsection{Stratégie de déploiement et gestion des règles}

La mise en service est progressive. Suricata fonctionne d'abord en mode IDS, ce qui permet d'observer les alertes sans risque pour le trafic. Après réglage des faux positifs, certaines règles à haute confiance passent en mode IPS (action \texttt{drop}). Le jeu de règles combine ET Open et des règles locales couvrant les flux DICOM et HL7, les accès d'administration anormaux et les comportements de balayage ou de force brute. Chaque alerte est journalisée dans \texttt{eve.json}, conservée pour analyse et relayée vers Wazuh.

\subsection{Éléments complémentaires}

La segmentation par zones et la supervision par Wazuh complètent le système sans en être l'objet : la première donne à la sonde de la visibilité sur les flux inter-zones, la seconde facilite la lecture des alertes. Elles sont mises en place dans le laboratoire à ce titre, et leur généralisation au réseau de production relève des recommandations.

\addcontentsline{toc}{section}{Conclusion}
\section*{Conclusion}

La problématique reste fondée sur les besoins exprimés par le service et sur les menaces hospitalières documentées au chapitre~\ref{chap:1} ; le nouvel audit prévu permettra de la confirmer par des mesures. En l'état, la proposition d'architecture repose sur Suricata, placé sur OPNsense, inspectant les flux entre cinq zones, d'abord en IDS puis en IPS. Le chapitre suivant décrit le protocole de déploiement en laboratoire et la méthode d'évaluation prévue, dont l'exécution reste à conduire.

 
\chapter{Déploiement de Suricata, règles personnalisées et évaluation avant/après}\label{chap:3}
 \addcontentsline{toc}{section}{Introduction}
\section*{Introduction}

Ce chapitre décrit le déploiement de Suricata sur la plateforme OPNsense dans un environnement virtualisé, la rédaction de règles adaptées au contexte hospitalier, puis l'évaluation du système par des scénarios d'attaque rejoués avant et après sa mise en œuvre. Conformément à la règle suivie tout au long de ce travail, aucun résultat n'est présenté tant qu'il n'a pas été mesuré : les passages marqués en rouge sont à compléter avec les résultats réels, et chaque section indique les preuves à produire.

\section{Environnement de virtualisation}

\subsection{Hyperviseur et ressources}

L'ensemble des manipulations est réalisé dans un environnement de test distinct du réseau de production, afin d'éviter tout risque d'interruption de service. L'hyperviseur retenu est KVM, piloté par libvirt, sur un hôte Ubuntu Server \acompleter{version}, avec \acompleter{X} Go de RAM et \acompleter{X} cœurs. La topologie est administrée via \acompleter{GNS3 ou EVE-NG}. Le tableau~\ref{tab:vms} liste les machines du laboratoire.

\begin{table}[htbp]
\centering
\small
\begin{tabularx}{\textwidth}{lXl}
\hline
\textbf{Machine virtuelle} & \textbf{Rôle} & \textbf{Ressources} \\
\hline
OPNsense & pare-feu et sonde Suricata & \acompleter{vCPU / RAM} \\
Wazuh (tout-en-un) & visualisation des alertes & \acompleter{vCPU / RAM} \\
Contrôleur Active Directory simulé & annuaire, DNS & \acompleter{vCPU / RAM} \\
Poste Windows 10 & poste utilisateur & \acompleter{vCPU / RAM} \\
Serveur Linux & serveur applicatif, cible des tests & \acompleter{vCPU / RAM} \\
Poste Kali & machine d'attaque (tests) & \acompleter{vCPU / RAM} \\
\hline
\end{tabularx}
\caption{Machines virtuelles du laboratoire}
\label{tab:vms}
\end{table}

\subsection{Installation et vérification}

La virtualisation matérielle est vérifiée par la commande \texttt{egrep -c '(vmx|svm)' /proc/cpuinfo}, qui doit renvoyer une valeur non nulle ; les machines créées sont listées par \texttt{virsh list --all}. Un pont réseau relie l'hôte au réseau physique du laboratoire.

\preuves{résultat de \texttt{egrep -c '(vmx|svm)' /proc/cpuinfo} ; capture de \texttt{virsh list --all} ; schéma de la topologie ; configuration du pont réseau de l'hôte.}

\subsection{Limites méthodologiques}\label{sec:limites}

Travailler sur un environnement virtualisé constitue une limite assumée : les mesures de performance sont relatives et non représentatives d'une charge réelle de production, le nombre de cœurs partagés entre toutes les machines peut affecter la précision des mesures de latence, et l'édition communautaire des outils peut différer d'une édition professionnelle. La comparaison avant/après reste néanmoins valide, car tous les états sont mesurés dans des conditions matérielles identiques.

\section{Plateforme : zones et pare-feu}

Les cinq zones du tableau~\ref{tab:vlan} sont créées sur le commutateur virtuel et rattachées chacune à une interface ou sous-interface d'OPNsense. OPNsense \acompleter{version} est installé depuis l'image ISO, avec un partitionnement ZFS \acompleter{mode retenu}. Le pare-feu applique le refus par défaut ; seuls les flux du tableau~\ref{tab:flux} transitent entre zones. Le durcissement des commutateurs (limitation des adresses MAC par port, DHCP snooping, \emph{Dynamic ARP Inspection}) n'est applicable que si le commutateur virtuel le permet \acompleter{indiquer ce qui a été effectivement activé}. Cette plateforme est identique dans tous les états d'évaluation : elle donne à la sonde de la visibilité sur les flux inter-zones.

\preuves{interfaces OPNsense avec chaque VLAN ; export des règles de pare-feu par zone ; pour un flux autorisé et un flux interdit, capture du test et ligne de journal correspondante ; résultat de \texttt{ping} et \texttt{traceroute} entre deux VLAN.}

\section{Déploiement de Suricata en mode IDS}

\subsection{Installation et configuration}

Suricata est activé sur les interfaces à surveiller d'OPNsense, avec le jeu de règles ET Open téléchargé et activé, d'abord en mode IDS (journalisation seule)\cite{opnsensedoc,suricatadoc}. Les variables d'adresses, en particulier le réseau protégé et les groupes de serveurs, sont définies selon le plan d'adressage du tableau~\ref{tab:vlan}, car de leur exactitude dépend la pertinence des règles. Le journal \texttt{eve.json} est activé.

\preuves{activation de Suricata et interfaces surveillées ; jeu de règles ET Open chargé ; variables d'adresses ; capture d'un premier événement dans \texttt{eve.json}.}

\subsection{Réglage et gestion des faux positifs}

Un jeu de règles complet génère des alertes sur du trafic légitime. Une période d'observation en IDS, sur un trafic normal du laboratoire, permet d'identifier les règles bruyantes, puis de les désactiver, de les restreindre à certaines zones ou de supprimer les alertes pour des hôtes précis. Cette étape conditionne le passage en IPS : une règle qui déclenche des faux positifs ne doit pas bloquer.

\preuves{nombre d'alertes sur trafic normal avant et après réglage ; liste des règles désactivées ou restreintes, avec justification.}

\section{Règles de détection personnalisées}

Des règles locales complètent ET Open pour le contexte hospitalier. Les exemples suivants couvrent les flux DICOM et HL7, un accès SMB, la force brute sur SSH et le balayage de ports. Ce sont des propositions, à tester et à adapter avant d'être présentées comme validées ; les variables sont définies dans la configuration de Suricata.

\begin{lstlisting}[language={},breaklines=true,basicstyle=\ttfamily\scriptsize]
alert tcp !$ZONE_SOINS_BIOMED any -> $SERVEURS_IMAGERIE 104 (msg:"CNHU DICOM depuis zone non autorisee"; flow:to_server; sid:1000001; rev:1;)
alert tcp any any -> $SERVEURS 445 (msg:"CNHU SMB vers serveurs depuis zone non autorisee"; flow:to_server; sid:1000002; rev:1;)
alert tcp !$ZONE_SOINS any -> $SERVEURS_SIH 2575 (msg:"CNHU HL7 depuis zone non autorisee"; flow:to_server; sid:1000003; rev:1;)
alert tcp any any -> $SERVEURS 22 (msg:"CNHU SSH tentatives repetees"; flags:S; threshold:type both, track by_src, count 5, seconds 60; sid:1000004; rev:1;)
alert tcp any any -> $HOME_NET any (msg:"CNHU balayage de ports"; flags:S; threshold:type both, track by_src, count 20, seconds 5; sid:1000005; rev:1;)
\end{lstlisting}

Pour chaque règle, le mémoire présente son contenu, l'objectif de détection visé et le résultat du test qui la déclenche.

\preuves{extrait des règles locales effectivement chargées ; pour chaque règle, l'alerte générée par le scénario correspondant (extrait de \texttt{eve.json} ou capture).}

\section{Passage en mode IPS}

Le mode IPS place Suricata en coupure sur le flux. Il repose sur la compatibilité de la carte réseau virtuelle avec le mécanisme de capture utilisé, ce qui doit être vérifié sur l'environnement de test avant la bascule. Seules les règles à haute confiance, validées pendant la phase IDS, passent à l'action \texttt{drop}. La charge ajoutée par l'inspection en ligne est mesurée par la latence et le débit avant et après l'activation de l'IPS.

\preuves{capture du blocage effectif d'une attaque en mode IPS, avec la ligne de journal ; mesure de latence ou de débit avant et après activation de l'IPS.}

\section{Visualisation des alertes avec Wazuh}

Wazuh est déployé en configuration tout-en-un pour centraliser et visualiser les alertes de Suricata\cite{wazuhdoc}. Les journaux d'OPNsense et les événements Suricata lui sont transmis selon la méthode retenue (agent sur OPNsense ou syslog). Une règle de corrélation peut élever le niveau d'une alerte donnée, et une réponse active peut bloquer temporairement l'adresse source après des échecs d'authentification répétés. Ce volet est complémentaire : l'évaluation porte sur Suricata.

\preuves{tableau de bord Wazuh accessible ; alerte Suricata visible dans le tableau de bord ; blocage automatique déclenché, si mis en œuvre.}

\section{Scénarios de test et protocole}

Sept scénarios, rattachés au référentiel MITRE ATT\&CK\cite{attck}, sont proposés (tableau~\ref{tab:scenarios}). Ils couvrent la reconnaissance, l'accès initial et le mouvement latéral. Chacun est rejoué à l'identique dans les trois états, exclusivement dans le laboratoire.

\begin{table}[htbp]
\centering
\small
\begin{tabularx}{\textwidth}{llXl}
\hline
\textbf{Réf.} & \textbf{Technique} & \textbf{Scénario} & \textbf{Outil} \\
\hline
S1 & T1018 & découverte des hôtes distants (balayage ping) & Nmap \\
S2 & T1046 & découverte des services réseau & Nmap \\
S3 & T1110 & force brute sur un service d'administration (SSH) & Hydra \\
S4 & T1557.002 & empoisonnement de cache ARP & arpspoof \\
S5 & T1021.002 & accès SMB depuis une zone non autorisée & smbclient \\
S6 & T1021.001 & connexion RDP depuis la zone Invités & client RDP \\
S7 & -- & accès DICOM depuis une zone non autorisée (flux métier interdit) & netcat \\
\hline
\end{tabularx}
\caption{Scénarios d'attaque simulés}
\label{tab:scenarios}
\end{table}

Trois états sont comparés, avec la même topologie et les mêmes règles de pare-feu : (A) l'état « avant », sans détection ; (B) Suricata en mode IDS ; (C) Suricata en mode IPS. Seul le fonctionnement de Suricata varie. Les indicateurs mesurés sont : attaque réussie ou non, alerte générée oui ou non, blocage oui ou non, délai de détection en secondes (horodatage de l'alerte moins horodatage du début du test), nombre de faux positifs et impact sur la latence. Le taux de détection est le nombre de scénarios détectés divisé par le nombre de scénarios joués ; le taux de blocage se calcule de la même façon.

\section{Résultats de l'évaluation avant/après}

Le tableau~\ref{tab:resultats} reste à renseigner avec les mesures réelles ; aucune valeur n'est à y inscrire avant les tests.

\begin{table}[htbp]
\centering
\small
\begin{tabular}{lccccc}
\hline
 & \textbf{(A) Avant} & \multicolumn{2}{c}{\textbf{(B) IDS}} & \multicolumn{2}{c}{\textbf{(C) IPS}} \\
\textbf{Scénario} & Alerte & Alerte & Délai (s) & Alerte & Blocage \\
\hline
S1 & \amesurer & \amesurer & \amesurer & \amesurer & \amesurer \\
S2 & \amesurer & \amesurer & \amesurer & \amesurer & \amesurer \\
S3 & \amesurer & \amesurer & \amesurer & \amesurer & \amesurer \\
S4 & \amesurer & \amesurer & \amesurer & \amesurer & \amesurer \\
S5 & \amesurer & \amesurer & \amesurer & \amesurer & \amesurer \\
S6 & \amesurer & \amesurer & \amesurer & \amesurer & \amesurer \\
S7 & \amesurer & \amesurer & \amesurer & \amesurer & \amesurer \\
\hline
\end{tabular}
\caption{Résultats de l'évaluation : sans détection, IDS et IPS}
\label{tab:resultats}
\end{table}

\preuves{pour chaque scénario et chaque état : commande lancée, sortie, présence ou absence d'alerte ; pour les états IDS et IPS, capture de l'alerte avec son horodatage ; nombre de faux positifs relevés ; mesure de latence avec et sans IPS.}

\subsection{Interprétation et limites}

\acompleter{pour chaque indicateur, une ou deux phrases sur ce que le résultat démontre par rapport à l'objectif : détecter et bloquer les attaques en préservant la disponibilité}

Les mesures de délai doivent être interprétées avec prudence : réalisées sur un environnement virtualisé aux cœurs partagés, elles ne reflètent pas nécessairement les performances de l'infrastructure de production du CNHU-HKM, dimensionnée différemment (voir la section~\ref{sec:limites}). Les résultats portent par ailleurs sur sept scénarios et sur la détection par signatures ; ils ne couvrent ni le trafic chiffré ni les menaces inédites.

\section{Difficultés rencontrées}

\acompleter{pour chaque difficulté du journal de bord : nature (message d'erreur exact), cause, solution retenue, temps perdu ; classer par composant (virtualisation, OPNsense, Suricata, Wazuh)}

\section{Recommandations pour un déploiement au CNHU-HKM}

\subsection{Préalables organisationnels}

Avant toute intervention sur la production, il convient d'obtenir la validation de la direction, de définir une fenêtre de maintenance et d'informer les services concernés, en particulier le service biomédical, dont les équipements ne tolèrent pas toujours une coupure ni un changement d'adressage.

\subsection{Déploiement progressif}

Un déploiement par étapes est recommandé : d'abord une sonde Suricata en mode passif, alimentée par un port miroir, sans effet sur le trafic ; puis une période de réglage des faux positifs ; ensuite un service pilote avec Suricata en mode IDS sur le pare-feu ; enfin, seulement après examen des faux positifs, le passage de règles à haute confiance en mode IPS. Cette prudence découle de la criticité des services de soins. La segmentation du réseau et la supervision des alertes avec Wazuh viennent compléter le dispositif.

\subsection{Moyens nécessaires}

Les besoins en matériel et en temps sont \acompleter{à estimer avec le service, en restant réaliste par rapport aux ressources observées pendant le stage}. Côté compétences, l'unité a elle-même identifié le besoin de développer l'expertise interne en cybersécurité, par des formations certifiantes disponibles au Bénin et des partenariats avec l'agence nationale chargée de la sécurité des systèmes d'information\cite{kaho2026} ; une prise en main de Suricata et d'OPNsense par les administrateurs en est un préalable. Les mesures rejoignent les recommandations de l'APDP sur les accès, la sécurisation des réseaux et la sauvegarde\cite{apdp2024}.

\addcontentsline{toc}{section}{Conclusion}
\section*{Conclusion}

Ce chapitre a décrit le déploiement de Suricata en modes IDS puis IPS sur OPNsense, les règles adaptées au contexte hospitalier, le protocole d'évaluation et les recommandations de déploiement. \acompleter{une phrase sur les résultats obtenus, à rédiger après les mesures}

 
% \include{perspectives}
%%conclusion
\conclusion

Ce mémoire s'est donné pour objectif de mettre en place, au CNHU-HKM, un système de détection et de prévention d'intrusion capable de repérer et de bloquer les attaques visant le réseau, sans compromettre la disponibilité des services de soins. L'audit a montré l'absence de segmentation visible et de détection d'intrusion sur le segment analysé ; après comparaison des solutions, Suricata a été déployé sur OPNsense, d'abord en mode IDS puis en mode IPS, avec des règles adaptées aux flux hospitaliers, et évalué en laboratoire par des scénarios rattachés à MITRE ATT\&CK. \acompleter{résumé des résultats du chapitre~\ref{chap:3} et de ce qu'ils démontrent par rapport à la problématique}

Ce travail apporte au CNHU-HKM un état des lieux documenté, un système de détection testé et des recommandations de déploiement progressif. Ses limites doivent être assumées : l'évaluation s'est faite en environnement virtualisé, avec un nombre limité de scénarios, et le périmètre ne couvre pas l'ensemble de l'infrastructure réelle.

Plusieurs perspectives se dégagent : segmentation complète du réseau par zones, supervision centralisée des alertes avec Wazuh, haute disponibilité du pare-feu, contrôle d'accès réseau (802.1X) et automatisation avec Ansible.

% 
\lhead[]{} \rhead[]{} \chead[]{}

%%biblio
\addcontentsline{toc}{chapter}{Bibliographie}
\bibliographystyle{abbrvnat}
\bibliography{biblio}


\chapter*{Annexes}\addcontentsline{toc}{chapter}{Annexes}

Les annexes regroupent les éléments justificatifs du mémoire. Celles qui correspondent à des documents ou à des mesures non encore disponibles portent un marquage explicite : elles ne contiennent aucune donnée reconstituée.

\section*{Annexe A -- Journal de bord du stage}
\addcontentsline{toc}{section}{Annexe A -- Journal de bord du stage}

Le journal de bord atteste du déroulement réel du stage. Seul son squelette est fourni ici (tableau~\ref{tab:journal}) ; les lignes doivent être complétées à partir des notes tenues pendant la période de stage, sans reconstruction a posteriori.

\begin{table}[htbp]
\centering
\small
\begin{tabularx}{\textwidth}{llXl}
\hline
\textbf{Date} & \textbf{Phase} & \textbf{Activité réalisée} & \textbf{Problème / solution} \\
\hline
\acompleter{ligne 1} & & & \\
\acompleter{ligne 2} & & & \\
\acompleter{lignes suivantes} & & & \\
\hline
\end{tabularx}
\caption{Trame du journal de bord, à compléter avec les entrées effectives}
\label{tab:journal}
\end{table}

\section*{Annexe B -- Fiche de durcissement d'OPNsense (proposition)}
\addcontentsline{toc}{section}{Annexe B -- Fiche de durcissement d'OPNsense}

Cette fiche est une proposition méthodologique ; les points cochés devront correspondre à une configuration effectivement appliquée et vérifiée sur l'installation du laboratoire.

\begin{itemize}
\item politique de filtrage en refus par défaut, seuls les flux du tableau de référence étant ouverts ;
\item restriction de l'interface d'administration au seul VLAN Administration, avec interdiction d'accès depuis les zones Invités et depuis l'interface WAN ;
\item comptes d'administration nominatifs, mots de passe robustes et, si l'édition le permet, authentification multifacteur ;
\item planification des mises à jour de sécurité et figement de la version utilisée, à reporter dans le tableau des versions du chapitre~\ref{chap:3} ;
\item désactivation des services non utilisés ;
\item activation de la journalisation et redirection des journaux par syslog vers le serveur Wazuh ;
\item sauvegarde de la configuration avant chaque modification et conservation chiffrée des exports ;
\item promotion d'une règle Suricata en action \texttt{drop} uniquement après la période d'observation en IDS.
\end{itemize}

\preuves{export de la configuration effective (\texttt{config.xml}) et captures d'écran des écrans correspondant à chaque point appliqué.}

\section*{Annexe C -- Configurations du laboratoire}
\addcontentsline{toc}{section}{Annexe C -- Configurations du laboratoire}

Les fichiers de configuration suivants doivent être joints dans leur version effective, accompagnés de leur date d'export :

\begin{itemize}
\item configuration complète d'OPNsense (export \texttt{config.xml}) : interfaces, VLAN, règles de filtrage ;
\item fichier de configuration de Suricata : variables d'adresses (cf. tableau des variables du chapitre~\ref{chap:3}), interfaces surveillées, activation de \texttt{eve.json} ;
\item fichier des règles locales, reproduit en annexe~D ;
\item configuration du pont réseau de l'hôte de virtualisation et caractéristiques des machines virtuelles.
\end{itemize}

\acompleter{insérer ici les exports réels ; aucun extrait de configuration n'est fabriqué en l'absence du laboratoire exécuté.}

\section*{Annexe D -- Jeu de règles locales Suricata}
\addcontentsline{toc}{section}{Annexe D -- Jeu de règles locales Suricata}

Le fichier \texttt{local.rules} proposé reprend les cinq règles présentées et commentées au chapitre~\ref{chap:3}. Il n'est réputé chargé qu'après un test de configuration réussi et le déclenchement vérifié de chaque règle.

\begin{lstlisting}[language={},breaklines=true,basicstyle=\ttfamily\scriptsize]
# R1 - trafic DICOM (104) vers l'imagerie depuis une zone non autorisee
alert tcp !$ZONE_SOINS_BIOMED any -> $SERVEURS_IMAGERIE 104 (msg:"CNHU DICOM depuis zone non autorisee"; flow:to_server; classtype:policy-violation; sid:1000001; rev:2;)
# R2 - acces SMB (445) vers les serveurs
alert tcp any any -> $SERVEURS 445 (msg:"CNHU SMB vers serveurs"; flow:to_server; classtype:policy-violation; sid:1000002; rev:2;)
# R3 - trafic HL7/MLLP (2575) vers le SIH depuis une zone non autorisee
alert tcp !$ZONE_SOINS any -> $SERVEURS_SIH 2575 (msg:"CNHU HL7 depuis zone non autorisee"; flow:to_server; classtype:policy-violation; sid:1000003; rev:2;)
# R4 - force brute SSH : 5 SYN d'une meme source en 60 s
alert tcp any any -> $SERVEURS 22 (msg:"CNHU SSH tentatives repetees"; flow:to_server; flags:S; classtype:attempted-admin; detection_filter:track by_src, count 5, seconds 60; sid:1000004; rev:2;)
# R5 - balayage : 20 SYN d'une meme source en 5 s
alert tcp any any -> $HOME_NET any (msg:"CNHU balayage de ports"; flags:S; classtype:attempted-recon; detection_filter:track by_src, count 20, seconds 5; sid:1000005; rev:2;)
\end{lstlisting}

Les commandes de contrôle (\texttt{suricata -T}, rejeu de capture et interrogation de \texttt{eve.json}) sont celles du chapitre~\ref{chap:3} ; leurs sorties réelles doivent être archivées en annexe~E.

\section*{Annexe E -- Résultats bruts des scans et des tests}
\addcontentsline{toc}{section}{Annexe E -- Résultats bruts des scans et des tests}

Cette annexe accueille les sorties non modifiées des commandes exécutées. Pour l'audit réseau, doivent être archivées les sorties complètes de la découverte d'hôtes (\texttt{nmap -sn}) et du balayage ciblé avec détection de version (\texttt{nmap -sV}) sur les sept ports, dans les formes exactes utilisées le 9~octobre. Pour l'évaluation, le tableau~\ref{tab:preuves} précise, pour chaque scénario, les éléments à joindre.

\begin{table}[htbp]
\centering
\small
\begin{tabularx}{\textwidth}{clX}
\hline
\textbf{Scénario} & \textbf{États} & \textbf{Éléments à archiver} \\
\hline
S1 à S7 & A, B, C & commande lancée et sortie complète ; horodatage de début ; extrait \texttt{eve.json} (états B et C) ; preuve de blocage (état C) \\
\hline
\end{tabularx}
\caption{Manifestation des preuves brutes à recueillir par scénario}
\label{tab:preuves}
\end{table}

\acompleter{déposer les fichiers de sortie réels ; en l'absence d'exécution, cette annexe reste vide et le chapitre~\ref{chap:3} conserve ses emplacements de mesure.}

\section*{Annexe F -- Matrice de flux complète}
\addcontentsline{toc}{section}{Annexe F -- Matrice de flux complète}

La matrice de flux est une proposition de politique, construite à partir des usages décrits avec le service ; elle devra être validée par l'UGRIM avant tout déploiement. Sa version de synthèse constitue le tableau de flux du chapitre~\ref{chap:2} ; les ports applicatifs du SIH et des équipements biomédicaux restent à confirmer.

\acompleter{étendre la matrice avec l'ensemble des services effectivement recensés (ports, sens, journalisation) après validation par le service.}

\section*{Annexe G -- Schémas réseau avant et après}
\addcontentsline{toc}{section}{Annexe G -- Schémas réseau avant et après}

Les figures du chapitre~\ref{chap:2} (architecture existante et architecture cible) sont actuellement des schémas TikZ simplifiés. Les schémas d'origine du service et les schémas cibles au format vectoriel (PDF ou SVG, par exemple exportés de draw.io) doivent être joints ici, après obtention de l'autorisation de reproduction du document interne.

\acompleter{insérer les schémas validés par le service informatique.}

\section*{Annexe H -- Autorisation de stage et cadrage du sujet}
\addcontentsline{toc}{section}{Annexe H -- Autorisation de stage et cadrage du sujet}

Cette annexe doit contenir le document signé par le tuteur de stage fixant : le périmètre autorisé pour la consultation de la documentation réseau ; l'autorisation des opérations de découverte (balayage d'hôtes et de ports) ; l'interdiction des tests d'intrusion actifs sur le réseau de production sans accord explicite ; l'engagement de ne pas consulter de données patients ; les règles d'anonymisation des informations sensibles.

\acompleter{insérer l'autorisation signée ; le mémoire ne peut présenter les manipulations comme autorisées sans ce document.}


\end{document}
